% ======================================================================
% 2.2 Цилиндрические координаты. Ответственный: D — Филиппов Андрей.
%
% ЭТА ЖЕ ЗАДАЧА разобрана в report_template.tex (строки 92–200) —
% главный ориентир оформления: блоки, таблица, рисунки, вывод.
%
% Что сдать (main.tex: условие 490–526, порядок 527–553, пример 670–706;
% требования к отчёту 1862–1870):
%   1. Изображение половины цилиндра: радиус, высота, граница y=0.
%   2. Обоснование пределов r, theta, z; замена координат; якобиан
%      вычислением определителя матрицы частных производных.
%   3. Повторный интеграл от g(r,theta,z) с явными пределами; аналитическое решение.
%   4. Метод средних прямоугольников при n=8,16,32,64;
%      шаги h_r=R/n, h_theta=pi/n, h_z=H/n.
%   5. Проверка объёма: интегрировать r (якобиан), сравнить с V=pi*R^2*H/2.
%   6. Таблица значений и ошибок, график ошибки от h=1/n, вывод о сходимости.
%
% Параметры варианта 3 (приложение А): R=1, H=3.
% Код: solution.py (блок D), запуск python experiments.py 2.2.
% Рисунки: figures/2-2-*.pdf.
% ======================================================================

\subsection{Цилиндрические координаты}

\taskhead{Исходные данные и аналитическое решение}
Для варианта 3 заданы $R=1$, $H=3$. Рассматривается интеграл
\[
I_{\mathrm{ц}}=\iiint_{\Omega_{\mathrm{ц}}}(1+x^2+y^2+z)\,dx\,dy\,dz,
\qquad
\Omega_{\mathrm{ц}}=\{x^2+y^2\leq R^2,\ y\geq0,\ 0\leq z\leq H\}.
\]
\todo{D: описать область; обосновать пределы $r,\theta,z$; выписать замену координат; вычислить якобиан через определитель матрицы частных производных; записать функцию $g(r,\theta,z)$ с множителем якобиана и повторный интеграл; вычислить аналитически (пример: main.tex, строки 670–706; оформление: \texttt{report\_template.tex}, строки 95–131).}

\taskhead{Программный результат и сравнение}
\todo{D: листинг вызова своей функции из solution.py (образец: \texttt{report\_template.tex}, строки 138–149); указать шаги $h_r=R/n$, $h_\theta=\pi/n$, $h_z=H/n$.}
Аналитическое значение $I_{\mathrm{ц}}$: \todo{D: подставить}.
\begin{table}[htbp]
\centering\small
\caption{Сравнение численного интеграла с аналитическим значением $I_{\mathrm{ц}}$}
\label{tab:2-2}
\begin{tabular}{|r|r|r|r|r|}\hline
$n$ & $h$ & $I_n$ & $\Delta_n$ & $\delta_n$\\\hline
8 & — & — & — & —\\
16 & — & — & — & —\\
32 & — & — & — & —\\
64 & — & — & — & —\\
\hline\end{tabular}
\end{table}

\taskhead{Проверка объёма}
\todo{D: заменить исходную функцию на 1 (в новых координатах интегрировать $r$, а не 1); вычислить на всех сетках и сравнить с $V=\pi R^2H/2$; пояснить результат (образец: \texttt{report\_template.tex}, строки 173–176).}

\taskhead{Графическая проверка}
\taskfigure[0.6]{figures/2-2-region.pdf}{Половина цилиндра при $R=1$, $H=3$ (равный масштаб осей)}{fig:2-2-region}
\taskfigure[0.75]{figures/2-2-error.pdf}{Абсолютная ошибка от нормированного шага $h=1/n$ (логарифмические оси)}{fig:2-2-error}
\todo{D: пояснить график.}

\taskhead{Вывод}
\todo{D: вывод о сходимости; что подтверждает проверка объёма, а что — нет (образец: \texttt{report\_template.tex}, строки 195–200).}
